\FloatBarrier
\section{Correlative microscopy}

Correlative microscopy is a generic term describing the combination of multiple microscopical methods.
In the field of biology it is usually understood as a combination of light/fluorescence and electron microscopy \cite{mironov.2013,loussertfonta.2015}.
Nonetheless, it can be applied to most microscopic and tomographic techniques, als long as the measurement is done at exactly the same specimen with sufficient overlap \cite{loussertfonta.2015,daly.2017,mitchell.2019,mitchell.2021}.
Ideally it is done within the same device, though this is not essential.

Advantages of the combination of multiple techniques are
\begin{itemize}
    \item the combination of different information of one object/phase, like:
    \begin{itemize}
        \item trace element concentrations (see Section~\ref{sec:corr_laicpms} and \cite{kleiner.2022b}),
        \item crystallographic information \cite{rossler.2022},
        \item mechanical information (hardness) \cite{chen.2010,li.2025}, % was von Feng - noch nicht veröffentlicht!
        \item ...
    \end{itemize}
    \item information of a specimen across multiple scales \cite{daly.2017,mitchell.2019,mitchell.2021,kleiner.2022a},
    \item The resolution differences can improve the segmentation results of lower resolution methods, enabling new insights \cite{kleiner.2022b}.
\end{itemize}

For example, the combination of \ac{SEM}-\ac{BSE}, \ac{EDX} and \ac{EBSD} allows to identify phases chemically and distiguish phase borders even between chemically identical phases, but with a different crystal structure or differing crystal orientations.

\subsection{Image alignment}

Especially, if the images stem from different devices with vastly different modes, the image alignment is the most challenging part.
However, if the resulting images share enough features, they still can be aligned.
This can be done using
\begin{itemize}
    \item a fully manual alignment by rotating and transforming the images,
    \item {\color{red} semi-manual method by identifying at least 4 matching point-pairs between both images \cite{},}
    \item automatic methods from the field of computer vision/photogrammetry like \ac{SIFT} \cite{lowe.2004} or machine learning supported algorithms.
\end{itemize}

Usually, for few or single datasets this is done manually.
However, this can become tedious time-consuming and in many cases it is not very exact.

The main issues for combination of multi-device measurements are
\begin{enumerate}
    \item distortions,
    \item pixel aspect ratio,
    \item significant resolution differences,
    \item not enough matching features in both images,
    \item and re-identifying the measurement positions in the first place.
\end{enumerate}

To simplify the realignment, markers (called fiducial), can be used, to align different data sources to each other.
However, this is not always possible or expensive.

% \cite{miliszkiewicz.2015} Current approaches to calibration of LA-ICP-MS analysis